\subsection{Lower homotopy groups of spheres}

The same choice of a direction avoiding finitely many linear subspaces
contracts a family of loops with several parameters. Use the based
cubical description: a representative in degree $k\ge1$ is a uniformly
presented map on $I^k=[-1,1]^k$ whose restriction to the boundary is one
specified named point $a$. Homotopies relative to the boundary have value
$a$ there at every time.

\begin{theorem}[Vanishing below the dimension of the sphere]\label{idx-higher10}
For $1\le k<n$, every such representative on the cubical sphere
$S^n_\infty=\{x\in[-1,1]^{n+1}:\|x\|_\infty=1\}$ has a uniformly
presented nullhomotopy relative to its boundary. Hence
$\pi_k(S^n,a)=0$ for these representatives, also in the round-sphere
presentation. Every two presented points of $S^n$ with $n\ge1$ can be
joined by a uniformly presented path. The constructions use primitive
recursion in the dimensions, mesh size and rational data, for each fixed
input evaluator and modulus.
\end{theorem}
\begin{proof}
Write $f$ for the input map, with a modulus $\mu$ on names in the maximum
metric. Take $m=2^{\mu(5)+1}$ and subdivide each parameter axis into $m$
pieces of length $l=2/m$. Two points of a mesh cube have images within
$1/32$. At each grid vertex $v$ take a rational vector $p_v$ within $1/32$
of $f(v)$, clipping to $[-1,1]^{n+1}$ if necessary. On \emph{every}
boundary vertex use the same reading of $a$. This makes the subsequent
polygonal approximation exactly constant on each boundary face, including
where different faces meet.

Here are the triangulation and its estimates. Within each grid cube write
the local coordinates as $t_i\in[0,1]$. Sort them as
$t_{i_1}\ge\cdots\ge t_{i_k}$, breaking ties by label. The simplex
vertices $v_0,\ldots,v_k$ start at the lower corner and successively
increase coordinates $i_1,\ldots,i_k$ by one grid step. For the affine
interpolant $P$ take the corresponding convex combination of the $p_{v_j}$
with weights
\[
 1-t_{i_1},\quad t_{i_1}-t_{i_2},\quad\ldots,\quad
 t_{i_{k-1}}-t_{i_k},\quad t_{i_k}.
\]
Exchanging equal coordinates changes only a zero-weight intermediate
vertex. On a shared cube face the vertices off the face have weight zero,
and the remaining coordinates have the same sorted order on either side.
The formulas consequently agree on all overlaps. On any one simplex,
$P=p_{v_0}+\sum_{j=1}^k t_{i_j}(p_{v_j}-p_{v_{j-1}})$, so its
Lipschitz constant is at most $2k/l=km$. This bound is global: for a
rational segment, split it where it crosses the finitely many grid and
coordinate-order hyperplanes. An equation holding on the entire segment
requires no cut. The maximum-norm lengths of successive subsegments sum
to the original length since all run along the same segment. This proves
the rational bound, whence extension gives a PR evaluator on names.
There are $M=m^k k!$ simplices in the enumeration (grid indices and
permutations). At rational parameters each vertex value used differs from
$f$ at the parameter by at most $1/16$; convexity and approximation give
$\|P-f\|_\infty\le1/16$ throughout. The first homotopy
$\rho((1-s)f+sP)$ has denominator at least $15/16$. At boundary
parameters its value depends only on $s$.

The $k+1$ rational vertex readings of each simplex admit a common
annihilating nonzero rational covector $c$ supported on the first $k+2$
coordinates. Indeed there are $k+1$ homogeneous row equations and $k+2$
unknowns. Row reduction uses at most $k+1$ pivots; choose a free coordinate,
set it to $1$ and the other free coordinates to zero, then back-substitute.
For $j\ge0$ put $u(j)=(1,j,\ldots,j^{k+1},0,\ldots,0)$ in
$\mathbb Q^{n+1}$. Each polynomial $c\cdot u(j)$ is nonzero as a
polynomial and excludes at most $k+1$ candidate integers. Indeed division
at a root $r$ gives the factor $T-r$ by coefficient subtraction. Every
other distinct root remains a root of the quotient. Repeating this on a
list of distinct roots more numerous than the degree would give a root of
a nonzero constant. The proof just uses rational operations on finite
lists. In the table counting excluded pairs (simplex, candidate), the
total therefore is at most $(k+1)M$. At least one of
$0,\ldots,(k+1)M$ has no exclusion. For its first such $j$ put
$w=u(j)/\|u(j)\|_\infty$. Its norm is one.

The second homotopy is $\rho((1-s)P+sw)$. For $s\le1/8$ the
denominator exceeds $1/2$ by $\|P\|_\infty\ge15/16$. On a simplex with
covector $c$ and $s\ge1/8$, pairing with $c$ gives the lower bound
$|c\cdot w|/(8\|c\|_1)$ because $c\cdot P=0$. Thus the minimum $b$ of
$1/2$ and these $M$ positive rational bounds works globally. The
inequalities follow first on rational parameters and then on names.
Normalization with this denominator is $(2/b)$-Lipschitz, as in the loop
proof. In both homotopies a convex mixture of two bounded vectors changes
by at most the larger endpoint displacement plus twice the change of the
time parameter. This estimate, $km$, the supplied modulus and $2/b$
give a modulus for the concatenated homotopy $H$. At the time seam use
target error $2^{-p-1}$ per side and double the time scale. The boundary
track $h(s)=H(v,s)$ is common to all $v\in\partial I^k$, with
$h(0)=a$ and $h(1)=w$.

We now keep the boundary fixed. Put $r=\|u\|_\infty$ for a parameter
$u\in I^k$. First deform $f(u)$ to its compressed copy using
$f(u/\max(r,1-s/2))$. The denominator is at least $1/2$, and at $r=1$
the input stays on the boundary. Next, at time $s$, use
\[
 \begin{cases}
  H(2u,s),&r\le1/2,\\
  h(2s(1-r)),&r\ge1/2.
 \end{cases}
\]
At the seam the values agree because $2u$ is a boundary point there.
At time zero this is the compressed copy, and at $r=1$ it equals $a$.
At time one the formula is $h(\min(1,2(1-r)))$ everywhere. Contract
it to $a$ with $h((1-s)\min(1,2(1-r)))$. The formula with cases is
evaluated first on rational parameters. The inputs to $H,h$ on the two
pieces vary by at most four times the domain displacement, and a segment
with one endpoint on either side can be split at its rational intersection
with $r=1/2$. Thus a common precision request on the two pieces, tightened
by one output bit and two input bits, yields a global modulus. The
case-free initial and final formulas have bounded denominators and
Lipschitz reparameterizations. Extension gives the three homotopies on
names, and their concatenation is relative to the boundary. The radial
homeomorphism transports the result to the round sphere.

For path connectedness take the two specified points as separate vertices,
with one rational reading close to each. Apply the same argument to these
two zero-dimensional simplices: a covector in the first two coordinates
for each single reading excludes at most one integer, so one of $0,1,2$
works for both. The two normalized interpolations join each original
point to the same $w$; reverse the second path to reach the second point.
\end{proof}

For Lean the new finite indices are a base-$m$ tuple and a permutation of
$k$ coordinate labels. Permutations may be obtained by filtering the
length-$k$ tuples with entries below $k$ for absence of repeats; for a simple
upper bound the search may use $m^k k^k$ in place of $M$ when $k\ge1$.
Prove the interpolation formula over rational inputs and its agreement
at ties before extending it. For the root bound keep a list of the distinct
roots already removed and the remaining coefficient vector; each synthetic
division shortens that vector by one and preserves vanishing at each of the
remaining listed roots. This is a bounded table invariant. The input for
the avoidance scan is the matrix of covectors and its output is the index
$j$ together with the rational $b$. These data supply all precision
requests in the three final homotopies.
